% !PRINTED 604
\weberpage{604}
%<<CONT
entsprechende Basisform $\alpha_1 t_1 + \alpha_2 t_2$, die mit $\varphi$ associirt ist, bestimmen. Wir haben nach \S.\,146, (3), wenn $\omega_1, \omega_2$ eine Basis von $\mathfrak{o}$ ist, die ganzen Zahlen $a_{1,1}, a_{1,2}, a_{2,2}$ so zu bestimmen, dass $a_{1,1}, a_{2,2}$ positiv und m\"oglichst klein sind, und dass
\[ \alpha_1 = a_{1,1}\omega_1, \quad \alpha_2 = a_{1,2}\omega_1 + a_{2,2}\omega_2 \]
durch $\varphi$ theilbar werden. Setzen wir also $(\omega_1, \omega_2) = (1, \Theta)$, so wird

\weq{1}{\alpha_1 = a_{1,1}, \quad \alpha_2 = a_{1,2} + a_{2,2}\Theta.}

Es ist also $a_{1,1}$ die kleinste positive durch $\varphi$ theilbare rationale Zahl, und $a_{2,2}$ ist die kleinste positive Zahl, f\"ur die $a_{2,2}\Theta$ nach dem Modul $\varphi$ mit einer rationalen Zahl congruent wird.

Wir wollen dies auf den Fall anwenden, dass $\varphi$ ein Primfunctional $\pi$ ist. Ein solches kann hier nur vom ersten oder vom zweiten Grade sein, je nachdem seine absolute Norm gleich $p$ oder $=p^2$ ist. Wir unterscheiden also die beiden F\"alle:

1) \quad $N_a(\pi) = p$

oder

2) \quad $N_a(\pi) = p^2.$

Da hier $N(p) = p^2$ ist, so ist im ersten Falle $p$ in zwei Primfactoren zerlegbar, im zweiten Falle enth\"alt $p$ nur einen Primfactor, d. h. $p$ ist im K\"orper $\Omega$ selbst noch als Primzahl zu betrachten.

Bestimmen wir nun f\"ur $\varphi = \pi$ die Basis (1), so ergiebt sich zun\"achst $a_{1,1} = p$, und $a_{2,2}$ ist die kleinste positive ganze rationale Zahl, f\"ur die das Product $a_{2,2}\Theta$ modulo $\pi$ mit einer rationalen Zahl congruent wird. Daf\"ur ist aber nach \S.\,150, 3 die nothwendige und hinreichende Bedingung
\[ (a_{2,2}\Theta)^p \equiv a_{2,2}\Theta \pmod{\pi}, \]
oder, da nach dem \gesperrt{Fermat}'schen Satze $a_{2,2}^p \equiv a_{2,2}$ ist,

\weq{2}{a_{2,2}(\Theta^p - \Theta) \equiv 0 \pmod{\pi}.}

Es ist also entweder $\Theta^p - \Theta \equiv 0$, und dann ist $a_{2,2} = 1$, oder $\Theta^p - \Theta$ nicht $\equiv 0$, und dann ist $a_{2,2} = p$.

Im ersten Falle ist $\Theta$ und damit jede Zahl $\omega$ mit einer rationalen Zahl congruent. Es giebt also nicht mehr als $p$ nach dem Modul $\pi$ incongruente Zahlen, und folglich ist (\S.\,148, 3.) $p^f = p$, $f = 1$, d. h. $p$ in zwei Primfactoren ersten Grades zerlegbar.
